\footnotesize
\noindent
% Definisi Environment Tabel
\newenvironment{tasktable}
{
    \begin{longtable}{
        | @{\hspace{0pt}} >{\centering\arraybackslash}p{1cm} @{\hspace{0pt}} 
        | @{\hspace{0pt}} >{\centering\arraybackslash}p{3.1cm} @{\hspace{0pt}} 
        | @{\hspace{0pt}} >{\centering\arraybackslash}p{3.1cm} @{\hspace{0pt}} 
        | @{\hspace{0pt}} >{\centering\arraybackslash}p{3.1cm} @{\hspace{0pt}} 
        | >{\raggedright\arraybackslash}p{3.1cm} |}
        \hline
        \textbf{No.} & \textbf{Jenis} & \textbf{Tanggal Pengerjaan} & \textbf{Requester} & \textbf{Detail} \\ \hline
        \endfirsthead
        
        \hline
        \textbf{No.} & \textbf{Jenis} & \textbf{Tanggal Pengerjaan} & \textbf{Requester} & \textbf{Detail} \\ \hline
        \endhead
}
{ 
    \end{longtable} 
}

% Command untuk input data sesuai gambar
% #1: No
% #2: Judul/ID Tugas (Baris panjang)
% #3: Jenis (IMPROVE, dll)
% #4: Tanggal
% #5: Requester
% #6: Detail (Lorem Ipsum)
\newcommand{\taskentry}[6]{
    \multirow{2}{*}{#1} & \multicolumn{4}{l|}{#2} \\ \cline{2-5}
    & #3 & #4 & #5 & #6 \\ \hline
}